\documentclass[10pt,a4paper]{amsart}
\usepackage[margin=19mm]{geometry}
\usepackage{amsmath,amssymb,mathtools}
\usepackage[scr=rsfs,cal=euler]{mathalfa}
\usepackage{xfrac}
\usepackage[unicode,hidelinks,bookmarks=false]{hyperref}
\newcommand{\ie}{i.e.\ }
\newcommand{\cf}{cf.\ }
\newcommand{\resp}{respectively\ }
\newcommand{\Subsection}{\subsection}
\providecommand{\nobreakdash}{\nobreak\hbox{-}}
\setlength{\emergencystretch}{2em}
\multlinegap0pt
\usepackage{bm}
\usepackage{bbm}
\usepackage{dsfont}
\usepackage{xspace}
\usepackage{cancel}
\usepackage{mathtools}
\usepackage{amsthm}
\makeatletter
\@ifundefined{theorem}{\newtheorem{theorem}{Theorem}}{}
\makeatother
\title[Algebraic Graph Theory]{Algebraic Graph Theory}
\author{M Reza Salarian}
\date{Source: arXiv:2604.20890v1 (CC BY 4.0)}
\begin{document}
\maketitle

\noindent\emph{Source and licence.} Algebraic Graph Theory. Version arXiv:2604.20890v1 (CC BY 4.0), distributed under the Creative Commons Attribution 4.0 International licence, as recorded in the arXiv licence metadata for this exact version and shown on the abstract page https://arxiv.org/abs/2604.20890v1. Full rights notice: licenses/arxiv-2604.20890-CC-BY-4.0.txt.\par\smallskip

\noindent\textit{Editorial scope.} Counted unit: the Theorem whose source label is \texttt{None} in the article (source0.tex lines 6768--6771, source label \texttt{None}), delivered together with the article's own complete original proof of it (source0.tex lines 6773--6871, one \texttt{proof} environment of 99 lines). No mathematics is written, repaired or re-derived: statement and proof are the article's own text line for line. Required context retained uncounted: none. Every definition, display and label of those lines is retained unchanged. Omitted from this package and not redistributed: 1--6767, 6772--6772, 6872--9060. Nothing inside a retained excerpt is omitted or altered: every retained body is delivered exactly as the article's lines are written, so no omission receipt of any kind accompanies this package. Citations: the retained excerpts carry no citation command at all, so no bibliography entry of the article is transcribed and no reference is redistributed. Reference adaptation: none was needed and none was made; every reference inside the delivered unit resolves inside it, so no unresolved reference or citation is produced. Printed slips kept literal: no printed word, symbol, hypothesis or number of the counted statement or of its proof was altered. Layout and class: the article's own class is \texttt{book} with packages including inputenc, fontenc, babel, caption, subcaption, tikz, graphicx, listings, xcolor, geometry, amsmath, amssymb, amsthm, amsfonts; the wrapper is the lane's pinned \texttt{amsart} editorial wrapper at 10pt on A4, which restates the theorem-like environments the retained text needs and adds the article's own macro definitions that the retained text needs (none), with extra wrapper packages (bm, bbm, dsfont, xspace, cancel, mathtools). The delivered numbering is the unit's own. Assets: no retained span contains a figure environment, a graphics-inclusion command, a PDF-page inclusion, a table or a TikZ picture; no image file is copied into this package, the package declares no asset dependency and no image file is redistributed with it. Licence: version arXiv:2604.20890v1 of this article is distributed under the Creative Commons Attribution 4.0 International licence, as recorded in the arXiv licence metadata for this exact version and shown on the abstract page https://arxiv.org/abs/2604.20890v1. A full rights notice is delivered at licenses/arxiv-2604.20890-CC-BY-4.0.txt. Characters: every retained excerpt is pure ASCII, so the delivered engine, XeLaTeX with its default Latin Modern text font, reports no missing character.
\begin{theorem}[Hoffman--Singleton (core eigenvalue argument)]
Let \(\Gamma\) be a Moore graph of diameter \(2\) and valency \(k\).
Then \(k\in\{2,3,7,57\}\).
\end{theorem}

\begin{proof}
Let \(\Gamma\) be \(k\)-regular of diameter \(2\) and attain the Moore
bound. For diameter \(2\) the Moore bound is
\[
n := |V(\Gamma)| = 1 + k + k(k-1) = k^2 + 1.
\]
Moreover \(\Gamma\) has girth \(2\cdot 2 + 1 = 5\), so there are no
triangles or 4-cycles. The diameter \(2\) condition implies every pair
of non-adjacent vertices is at distance \(2\); together with absence of
4-cycles, this forces that any two distinct non-adjacent vertices have
\emph{exactly one} common neighbour, while any two adjacent vertices
have \emph{zero} common neighbours (no triangles). Thus for the adjacency
matrix \(A\) of \(\Gamma\) we have the combinatorial identity
\[
(A^2)_{ij} =
\begin{cases}
k, & i=j,\\[3pt]
0, & i\ne j \text{ and } A_{ij}=1,\\[3pt]
1, & i\ne j \text{ and } A_{ij}=0.
\end{cases}
\]
Equivalently, as a matrix equation,
\begin{equation}\label{eq:poly}
A^2 + A - (k-1)I \;=\; J,
\end{equation}
where \(J\) denotes the all-ones matrix and \(I\) the identity.

We use (\ref{eq:poly}) to deduce the spectrum of \(A\). Let \(\mathbf{1}\)
be the all-ones vector. Since \(\Gamma\) is \(k\)-regular,
\(A\mathbf{1}=k\mathbf{1}\) and \(J=\mathbf{1}\mathbf{1}^T\) acts as
\(J\mathbf{1}=n\mathbf{1}\) and annihilates any vector orthogonal to
\(\mathbf{1}\).

Apply (\ref{eq:poly}) to an eigenvector \(x\) with eigenvalue \(\theta\).

- If \(x\) is proportional to \(\mathbf{1}\), we recover the trivial
  eigenpair \(\theta=k\) and (\ref{eq:poly}) gives
  \(k^2 + k - (k-1) = n\), which holds since \(n=k^2+1\).

- If \(x\) is orthogonal to \(\mathbf{1}\), then \(Jx=0\) and
  (\ref{eq:poly}) reduces to the quadratic equation
  \[
  \theta^2 + \theta - (k-1) = 0.
  \]
  Hence every eigenvalue other than \(k\) is a root of this quadratic.
  Denote the two roots by
  \[
  r = \frac{-1 + \sqrt{4k-3}}{2},\qquad s = \frac{-1 - \sqrt{4k-3}}{2}.
  \]
  Let \(m_r\) and \(m_s\) be their multiplicities. We have \(1+m_r+m_s=n\).

We now use standard spectral identities (trace and sum of squares of
eigenvalues) to compute the multiplicities. The sum of all eigenvalues
equals \(\operatorname{tr}(A)=0\), hence
\[
k + m_r r + m_s s = 0.
\]
Also the sum of squares of eigenvalues equals \(\operatorname{tr}(A^2)=nk\),
so
\[
k^2 + m_r r^2 + m_s s^2 = nk = k^3 + k.
\]
Using \(m_r+m_s = n-1 = k^2\) these two linear equations in \(m_r,m_s\)
can be solved (or one may use standard manipulations) to obtain
\[
m_r = \frac{k(k-1) - (k+1)s}{r-s},\qquad
m_s = \frac{k(k-1) - (k+1)r}{s-r}.
\]
Since \(r-s = \sqrt{4k-3}\) is positive, the multiplicities \(m_r,m_s\)
are rational expressions in \(k\) and \(\sqrt{4k-3}\). Because
\(m_r,m_s\) must be nonnegative integers, strong arithmetic restrictions
arise. A (standard) simplification gives the following explicit closed
forms:
\[
m_r = \frac{k(k+1)(r+1)}{2r+1},\qquad
m_s = \frac{k(k+1)(s+1)}{2s+1},
\]
where \(2r+1=\sqrt{4k-3}\) and \(2s+1=-\sqrt{4k-3}\).

Set \(t := \sqrt{4k-3}\). Then \(t\) is a positive real number and
\(t^2 = 4k-3\). The above multiplicity formulas become rational
expressions in \(k\) and \(t\). Clearing denominators and using
integrality of \(m_r,m_s\) one obtains that \(t\) must be an odd integer
dividing \(k(k-1)\); in particular \(t\) is an odd positive integer. A
short divisibility and size argument (compare sizes of \(t\) and \(k\))
now forces that the only possible values of \(t\) are \(1,3,5,15\).
These correspond to \(k\) equal to
\[
k = \frac{t^2+3}{4} \in \{2,3,7,57\},
\]
respectively. Thus the degree \(k\) of a Moore graph of diameter \(2\)
must belong to \(\{2,3,7,57\}\).

(At this point one invokes classical facts: the cases \(k=2,3,7\) are
realized by the cycle \(C_5\), the Petersen graph and the
Hoffman--Singleton graph respectively; the case \(k=57\) is the only
remaining theoretical possibility and the existence of a Moore graph of
degree \(57\) is a famous deep problem; no other degrees occur.)
\end{proof}

\end{document}
